\documentclass{umk-matematika}

\begin{document}

\maketitlepage
    {Практическая работа №15}
    {Перпендикулярность плоскостей}
    {}

\section{Фундамент (1--4 балла)}

\textbf{Задача 1 (1 балл).} Сформулируйте признак перпендикулярности двух плоскостей.

\textbf{Задача 2 (2 балла).} Сформулируйте следствие из признака перпендикулярности плоскостей.

\textbf{Задача 3 (3 балла).} В кубе $ABCDA_1B_1C_1D_1$ перпендикулярны ли плоскости ABB₁A₁ и ABCD?

\textbf{Задача 4 (4 балла).} Приведите пример перпендикулярных плоскостей в строительстве.

\section{Стены (5--7 баллов)}

\textbf{Задача 5 (5 баллов).} Плоскость α проходит через катет BC прямоугольного треугольника ABC ($\angle C = 90^\circ$). Катет AC образует с плоскостью α угол 30$^\circ$. Найдите угол между плоскостью α и плоскостью треугольника.

\textbf{Задача 6 (6 баллов).} Через сторону AD прямоугольника ABCD проведена плоскость α так, что проекция BC на α равна 5 см. Найдите расстояние от BC до α, если AB = 12 см.

\textbf{Задача 7 (7 баллов).} Даны две перпендикулярные плоскости α и β. Прямая a лежит в плоскости α и параллельна линии пересечения плоскостей. Перпендикулярна ли прямая a плоскости β?

\section{Кровля (8--10 баллов)}

\textbf{Задача 8 (8 баллов).} В прямоугольном параллелепипеде $ABCDA_1B_1C_1D_1$ AB = 6 см, AD = 8 см, AA₁ = 10 см. Найдите угол между плоскостями AB₁C₁ и ABC.

\textbf{Задача 9 (9 баллов).} Плоскости α и β перпендикулярны. Точка A удалена от α на 6 см, от β на 8 см, от линии их пересечения на x см. Найдите x.

\textbf{Задача 10 (10 баллов).} Две стены здания перпендикулярны. На одной стене на высоте 3 м от пола закреплён кронштейн, на другой стене на высоте 4 м — другой кронштейн. Расстояние между проекциями кронштейнов на пол (вдоль линии пересечения стен) равно 12 м. Найдите расстояние между кронштейнами.

\newpage
\section*{Ответы}

\begin{enumerate}
  \item признак перпендикулярности плоскостей.
  \item следствие из признака.
  \item да.
  \item стена и пол.
  \item 60$^\circ$.
  \item $\sqrt{119}$ см.
  \item нет, прямая a параллельна плоскости β.
  \item 45$^\circ$.
  \item $x = 10$ см.
  \item 13 м.
\end{enumerate}

\end{document}